\chapter{Discrete Random Variables and Expectation}
\label{ch:drv}

Outcomes of experiments are often not numbers (a card, a sequence of tosses),
but we usually care about a numerical summary: the number of heads, the sum of
two dice, the waiting time for the first success. A \emph{random variable}
is the bridge between outcomes and numbers. This chapter defines random
variables on discrete sample spaces, introduces the probability mass function
and the cumulative distribution function, and develops the central quantity of
probability: the expectation.

\section{Random variables}

\begin{definition}[Random variable]
Let $(\Omega,\F,\Prob)$ be a probability space. A \emph{random variable} is a
function $X:\Omega\to\R$ such that for every real $x$, the set
$\{\omega:X(\omega)\le x\}$ is an event (that is, belongs to $\F$).
\end{definition}

A random variable is neither random nor a variable: it is a \emph{function} on
the sample space, and the randomness comes from the outcome $\omega$ being
random. The condition that $\{X\le x\}$ is an event guarantees that questions
such as ``what is the probability that $X\le x$?'' make sense. When $\Omega$ is
finite or countable and $\F=2^\Omega$, every function $X:\Omega\to\R$ satisfies
this condition automatically, so for now we can ignore it. Chapter~\ref{ch:measure}
revisits the condition.

\begin{definition}[Discrete random variable]
A random variable $X$ is \emph{discrete} if it takes values in a finite or
countable set $\{x_1,x_2,\dots\}$. The smallest such set (the set of values with
positive probability, together with possible limit points we shall not need) is
called the \emph{support} of $X$.
\end{definition}

We use the shorthand $\{X=x\}$ for the event $\{\omega\in\Omega:X(\omega)=x\}$ and
write $\Prob(X=x)$, $\Prob(X\in A)$, $\Prob(X\le x)$ accordingly.

\begin{example}
Roll two fair dice, $\Omega=\{1,\dots,6\}^2$. Let $S(i,j)=i+j$ and
$M(i,j)=\max(i,j)$. Both are random variables. The event $\{S=4\}$ is
$\{(1,3),(2,2),(3,1)\}$, with probability $3/36=1/12$.
\end{example}

\begin{definition}[Indicator random variable]
For an event $A$, the \emph{indicator} of $A$ is the random variable
$\indic_A(\omega)=1$ if $\omega\in A$ and $0$ otherwise.
\end{definition}

Indicators translate events into numbers. They satisfy
$\indic_{A\cap B}=\indic_A\indic_B$, $\indic_{A^c}=1-\indic_A$, and
$\indic_{A\cup B}=\indic_A+\indic_B-\indic_A\indic_B$. We will see that
expressing a count as a sum of indicators is one of the most powerful techniques
in the subject.

\section{Distribution: the mass function and the CDF}

\begin{definition}[Probability mass function]
The \emph{probability mass function} (PMF) of a discrete random variable $X$ is
$p_X(x)=\Prob(X=x)$, for $x\in\R$.
\end{definition}

\begin{proposition}\label{prop:pmf}
The PMF satisfies $p_X(x)\ge0$, $p_X(x)>0$ for at most countably many $x$, and
$\sum_xp_X(x)=1$. Moreover, for any $A\subseteq\R$,
$\Prob(X\in A)=\sum_{x\in A}p_X(x)$.
\end{proposition}
\begin{proof}
The events $\{X=x\}$, $x$ in the countable value set, form a partition of
$\Omega$. Apply countable additivity, first to the whole partition (giving the
sum $1$) and then to the sub-collection with $x\in A$.
\end{proof}

The PMF is the complete description of how $X$ is distributed; everything we
need about $X$ alone can be computed from it. Two different random variables
can have the same PMF.

\begin{nuance}{same distribution does not mean same random variable}
Toss a fair coin. Let $X=1$ if heads, $0$ if tails, and $Y=1-X$. Then $X$ and $Y$
each have PMF $p(0)=p(1)=1/2$, so they have the same \emph{distribution}, but
$X(\omega)\ne Y(\omega)$ for \emph{every} outcome $\omega$. We write
$X\overset{d}{=}Y$ for ``equal in distribution''. Equality in distribution says
nothing about how two random variables are related on the same sample space;
that information lives in the \emph{joint} distribution (Chapter~\ref{ch:joint}).
\end{nuance}

\begin{definition}[Cumulative distribution function]
The \emph{cumulative distribution function} (CDF) of a random variable $X$ is
$F_X(x)=\Prob(X\le x)$, for $x\in\R$.
\end{definition}

The CDF is defined for every random variable, discrete or not, which is why it is
the universal tool.

\begin{theorem}[Properties of the CDF]\label{thm:cdf}
Let $F=F_X$. Then
\begin{enumerate}[label=(\roman*)]
\item $F$ is non-decreasing;
\item $\lim_{x\to-\infty}F(x)=0$ and $\lim_{x\to+\infty}F(x)=1$;
\item $F$ is right-continuous: $\lim_{h\downarrow0}F(x+h)=F(x)$;
\item $\Prob(a<X\le b)=F(b)-F(a)$ for $a<b$;
\item $\Prob(X=x)=F(x)-F(x^-)$, where $F(x^-)=\lim_{h\downarrow0}F(x-h)$ is the left
limit; hence $F$ jumps at $x$ by exactly $p_X(x)$.
\end{enumerate}
\end{theorem}
\begin{proof}
(i) If $x\le y$ then $\{X\le x\}\subseteq\{X\le y\}$; apply monotonicity.

(ii) Let $x_n\to-\infty$ be any decreasing sequence. The events $\{X\le x_n\}$
decrease to $\bigcap_n\{X\le x_n\}=\emptyset$, since no real number $X(\omega)$ is
below every $x_n$. By continuity of probability (Theorem~\ref{thm:continuity}),
$F(x_n)\to\Prob(\emptyset)=0$. Because $F$ is monotone, the limit along any
sequence equals the limit along all of $\R$. For $+\infty$, the events
$\{X\le x_n\}$ with $x_n\uparrow\infty$ increase to $\Omega$, giving limit $1$.

(iii) Take $h_n\downarrow0$. The events $\{X\le x+h_n\}$ decrease to
$\{X\le x\}$ (since $X(\omega)\le x+h_n$ for all $n$ if and only if
$X(\omega)\le x$). Continuity from above gives the limit.

(iv) $\{X\le b\}$ is the disjoint union of $\{X\le a\}$ and $\{a<X\le b\}$.

(v) The events $\{X\le x-h_n\}$ increase to $\{X<x\}$. By continuity from below,
$F(x^-)=\Prob(X<x)$, and $\Prob(X=x)=\Prob(X\le x)-\Prob(X<x)$.
\end{proof}

\begin{nuance}{right-continuity, not left-continuity}
The CDF uses ``$\le$'', which makes it right-continuous; a CDF built with
``$<$'' would be left-continuous. When you evaluate a CDF at a jump, the value is
the \emph{upper} value of the jump. This is why $\Prob(a<X\le b)$, rather than
$\Prob(a\le X\le b)$, is the clean difference of CDF values.
\end{nuance}

\begin{example}[Max of two dice]\label{ex:maxdice}
Let $M=\max(i,j)$ for two fair dice. Then $\{M\le m\}$ is the event that both
rolls are at most $m$, with probability $m^2/36$. So $F_M(m)=m^2/36$ for
$m=1,\dots,6$, and by Theorem~\ref{thm:cdf}(v) the PMF is
$p_M(m)=(m^2-(m-1)^2)/36=(2m-1)/36$.
\end{example}

\subsection*{Functions of a random variable}

If $X$ is a random variable and $g:\R\to\R$ a function, then $Y=g(X)$ is a random
variable (for discrete $X$ always; the general condition is in
Chapter~\ref{ch:measure}). Its PMF is obtained by adding up the masses of all
$x$ that map to each $y$:
\[
p_Y(y)=\Prob(g(X)=y)=\sum_{x:\,g(x)=y}p_X(x).
\]

\begin{example}
Let $X$ be uniform on $\{-2,-1,0,1,2\}$ and $Y=X^2$. Then $Y\in\{0,1,4\}$ with
$p_Y(0)=1/5$, $p_Y(1)=2/5$, $p_Y(4)=2/5$. Note that $g$ is not injective, so the
masses of $x=\pm1$ merge.
\end{example}

\section{Expectation}

The expectation is the probability-weighted average of the values of $X$.

\begin{definition}[Expectation]
Let $X$ be a discrete random variable with PMF $p_X$. If
$\sum_x\abs{x}\,p_X(x)<\infty$, the \emph{expectation} (or mean) of $X$ is
\[
\Exp[X]=\sum_xx\,p_X(x).
\]
We then say $X$ is \emph{integrable}, written $\Exp\abs X<\infty$. If the
condition fails, the expectation is not defined (with one exception: if
$X\ge0$, we may write $\Exp[X]=+\infty$).
\end{definition}

\begin{nuance}{why absolute convergence is required}
Why insist on $\sum\abs x p(x)<\infty$ rather than just convergence of
$\sum xp(x)$? A conditionally convergent series can be rearranged to converge to
any number (Riemann's rearrangement theorem). The mean must not depend on the
order in which we list the values of $X$, and absolute convergence is exactly
what guarantees this. Concretely, take $\Prob(X=(-1)^{k}k)=c/k^2$ for
$k=1,2,\dots$ with $c=6/\pi^2$; then $\sum xp(x)=c\sum(-1)^k/k$ converges, but
$\sum\abs xp(x)=c\sum1/k=\infty$, and we declare $\Exp[X]$ undefined.
\end{nuance}

\begin{example}[The St.\ Petersburg game]\label{ex:stpete}
Toss a fair coin until the first head, at toss $K$, and receive $2^K$ rupees.
Then $\Prob(K=k)=2^{-k}$ and the payoff $X=2^K$ is non-negative with
\[
\Exp[X]=\sum_{k=1}^\infty2^k\cdot2^{-k}=\sum_{k=1}^\infty1=+\infty.
\]
The game has infinite expected payoff, yet nobody would pay a large sum to play.
Resolutions include limited bankroll or diminishing marginal utility (Bernoulli's
suggestion to use $\Exp[\log X]$). Note also that $X<\infty$ with probability
$1$: infinite expectation does not mean the payoff itself is infinite.
\end{example}

\begin{theorem}[Law of the unconscious statistician, LOTUS]\label{thm:lotus}
Let $X$ be discrete and $g:\R\to\R$. If $\sum_x\abs{g(x)}p_X(x)<\infty$, then
\[
\Exp[g(X)]=\sum_xg(x)\,p_X(x).
\]
\end{theorem}
\begin{proof}
Let $Y=g(X)$. Group the values $x$ according to $y=g(x)$:
\[
\sum_xg(x)p_X(x)=\sum_y\sum_{x:\,g(x)=y}y\,p_X(x)=\sum_yy\sum_{x:\,g(x)=y}p_X(x)
=\sum_yy\,p_Y(y)=\Exp[Y].
\]
Regrouping is legitimate because the series converges absolutely.
\end{proof}

LOTUS says we can compute $\Exp[g(X)]$ from the distribution of $X$ without first
finding the distribution of $g(X)$.

\begin{proposition}[Expectation over the sample space]\label{prop:expomega}
If $\Omega$ is countable with $p(\omega)=\Prob(\{\omega\})$, then
$\Exp[X]=\sum_{\omega}X(\omega)p(\omega)$, whenever
$\sum_\omega\abs{X(\omega)}p(\omega)<\infty$.
\end{proposition}
\begin{proof}
Group the outcomes by the value of $X$: $\sum_\omega X(\omega)p(\omega)=
\sum_xx\sum_{\omega:X(\omega)=x}p(\omega)=\sum_xx\,p_X(x)$.
\end{proof}

\begin{theorem}[Linearity of expectation]\label{thm:linear}
Let $X,Y$ be integrable random variables on a countable sample space and
$a,b\in\R$. Then $aX+bY$ is integrable and
\[
\Exp[aX+bY]=a\Exp[X]+b\Exp[Y].
\]
No assumption about the relationship between $X$ and $Y$ (in particular, no
independence) is needed.
\end{theorem}
\begin{proof}
By Proposition~\ref{prop:expomega},
\[
\Exp[aX+bY]=\sum_\omega(aX(\omega)+bY(\omega))p(\omega)
=a\sum_\omega X(\omega)p(\omega)+b\sum_\omega Y(\omega)p(\omega),
\]
where splitting the series is valid because both converge absolutely. Absolute
convergence of the left side follows from $\abs{aX+bY}\le\abs a\abs X+\abs b\abs Y$.
\end{proof}

(For random variables on a general probability space, linearity is proved in
Chapter~\ref{ch:measure}; the statement is the same.)

\begin{proposition}[Further properties]\label{prop:expprops}
\begin{enumerate}[label=(\roman*)]
\item If $X\ge0$ then $\Exp[X]\ge0$, with equality if and only if $\Prob(X=0)=1$.
\item (Monotonicity) If $X\le Y$ then $\Exp[X]\le\Exp[Y]$.
\item $\abs{\Exp[X]}\le\Exp\abs X$.
\item $\Exp[\indic_A]=\Prob(A)$.
\item (Tail-sum formula) If $X$ takes values in $\{0,1,2,\dots\}$, then
$\Exp[X]=\sum_{k=1}^\infty\Prob(X\ge k)$.
\end{enumerate}
\end{proposition}
\begin{proof}
(i) Each term $xp(x)$ with $x\ge0$ is non-negative; the sum is $0$ only if every
$x>0$ has $p(x)=0$. (ii) $Y-X\ge0$, so by (i) and linearity $\Exp[Y]-\Exp[X]\ge0$.
(iii) Apply (ii) to $-\abs X\le X\le\abs X$. (iv) $\indic_A$ takes value $1$ with
probability $\Prob(A)$ and $0$ otherwise.

(v) Write $X=\sum_{k=1}^\infty\indic_{\{X\ge k\}}$: for an outcome with $X=m$, the
indicators for $k=1,\dots,m$ equal $1$ and the rest $0$, so the sum is $m$. Taking
expectations term by term (justified because all terms are non-negative, so
we may interchange the double sum $\sum_k\sum_x$) gives the formula:
\[
\sum_{k\ge1}\Prob(X\ge k)=\sum_{k\ge1}\sum_{x\ge k}p(x)=\sum_{x\ge1}p(x)\sum_{k=1}^{x}1=\sum_xx\,p(x).\qedhere
\]
\end{proof}

\subsection*{The indicator method}

Linearity plus indicators gives a way to compute expectations of counts without
finding their distributions.

\begin{interview}{linearity of expectation with indicators}
To find $\Exp[N]$ for a count $N$: write $N=\sum_i\indic_{A_i}$ where each $A_i$ is
the event that item $i$ ``counts'', then
$\Exp[N]=\sum_i\Prob(A_i)$. Dependence among the $A_i$ is irrelevant. This
single idea solves a large fraction of expectation questions in interviews.
\end{interview}

\begin{example}[Fixed points of a random permutation]
Let $N$ be the number of fixed points of a uniformly random permutation of
$\{1,\dots,n\}$ (the matching problem, Example~\ref{ex:matching}). With
$A_i=\{\text{position }i\text{ is fixed}\}$, $\Prob(A_i)=(n-1)!/n!=1/n$, so
$\Exp[N]=n\cdot\frac1n=1$ for every $n$. Compare with the effort needed to find
the distribution of $N$; here we get the mean in one line, even though the $A_i$
are dependent.
\end{example}

\begin{example}[Records in a random sequence]\label{ex:records}
Let $x_1,\dots,x_n$ be a uniformly random ordering of $n$ distinct numbers. Say
position $k$ is a \emph{record} if $x_k$ exceeds all earlier values. Let $R_k$
indicate a record at position $k$. By symmetry, the largest among the first $k$
values is equally likely to be in any of the $k$ positions, so
$\Prob(R_k=1)=1/k$. Hence the expected number of records is the harmonic number
\[
\Exp\Bigl[\sum_kR_k\Bigr]=\sum_{k=1}^n\frac1k=H_n\approx\ln n+\gamma,\qquad
\gamma\approx0.5772 .
\]
For $n=100$ this is about $5.19$: a record-setting streak stays rare.
\end{example}

\begin{example}[Matching birthdays]
In a group of $n$ people with $365$ equally likely birthdays, the number of
\emph{pairs} sharing a birthday is $\sum_{i<j}\indic_{\{i,j\text{ match}\}}$, with
each indicator having expectation $1/365$. So the expected number of matching
pairs is $\binom n2/365$. For $n=23$ this is $253/365\approx0.69$, consistent with
the birthday calculation of Example~\ref{ex:birthday}.
\end{example}

\section{Variance}

The mean describes the centre of a distribution. The variance measures its
spread.

\begin{definition}[Variance]
If $\Exp[X^2]<\infty$, the \emph{variance} of $X$ is
$\Var(X)=\Exp\bigl[(X-\mu)^2\bigr]$, where $\mu=\Exp[X]$. The
\emph{standard deviation} is $\sigma_X=\sqrt{\Var(X)}$.
\end{definition}

(Finite second moment implies finite first moment, because
$\abs x\le(1+x^2)/2$.)

\begin{proposition}[Variance identities]\label{prop:var}
\begin{enumerate}[label=(\roman*)]
\item $\Var(X)=\Exp[X^2]-(\Exp[X])^2\ge0$.
\item $\Var(aX+b)=a^2\Var(X)$.
\item $\Var(X)=0$ if and only if $\Prob(X=\mu)=1$.
\item For any constant $c$, $\Exp[(X-c)^2]=\Var(X)+(\mu-c)^2$; hence the mean is
the constant that minimises the mean squared deviation.
\end{enumerate}
\end{proposition}
\begin{proof}
(i) Expand $(X-\mu)^2=X^2-2\mu X+\mu^2$ and use linearity:
$\Exp[X^2]-2\mu\cdot\mu+\mu^2=\Exp[X^2]-\mu^2$. Non-negativity holds because
$(X-\mu)^2\ge0$.

(ii) $aX+b$ has mean $a\mu+b$, so $(aX+b)-(a\mu+b)=a(X-\mu)$ and its square has
expectation $a^2\Var(X)$.

(iii) By Proposition~\ref{prop:expprops}(i), $\Exp[(X-\mu)^2]=0$ exactly when
$(X-\mu)^2=0$ almost surely.

(iv) $(X-c)^2=((X-\mu)+(\mu-c))^2=(X-\mu)^2+2(\mu-c)(X-\mu)+(\mu-c)^2$. Taking
expectations, the middle term vanishes since $\Exp[X-\mu]=0$.
\end{proof}

\begin{example}[Variance of an indicator]
For $X=\indic_A$ with $p=\Prob(A)$: $X^2=X$, so $\Var(X)=p-p^2=p(1-p)$. This is
maximised at $p=1/2$.
\end{example}

\begin{pitfall}{$\Var(X+Y)\ne\Var(X)+\Var(Y)$ in general}
Linearity of expectation holds always, but variance is additive only when the
variables are uncorrelated (in particular, when independent). We prove the exact
formula in Chapter~\ref{ch:joint}. Also, $\Var(-X)=\Var(X)$, not $-\Var(X)$.
\end{pitfall}

\begin{nuance}{the mean need not be a value $X$ can take}
A fair die has $\Exp[X]=3.5$, which is not a possible outcome. The mean is a
long-run average, not a ``typical'' value. The \emph{median} and \emph{mode} can
be more representative for skewed distributions, and for heavy-tailed variables
(Example~\ref{ex:stpete}) the mean may not exist at all.
\end{nuance}

\section*{Exercises}
\addcontentsline{toc}{section}{Exercises}

\begin{enumerate}
\item Find the PMF, mean and variance of the sum of two fair dice. (Use
linearity for the mean.)
\item Let $X$ have CDF $F(x)=0$ for $x<1$, $F(x)=0.3$ for $1\le x<3$,
$F(x)=0.7$ for $3\le x<5$, and $F(x)=1$ for $x\ge5$. Find the PMF and
$\Prob(2<X\le5)$.
\item Prove that if $X$ takes values in $\{0,1,2,\dots\}$ then
$\Exp[X^2]=\sum_{k\ge1}(2k-1)\Prob(X\ge k)$.
\item A deck of $n$ cards numbered $1,\dots,n$ is shuffled and the cards are
turned over one at a time. Find the expected number of times a card is
immediately followed by a card with a larger number.
\item (Expected number of distinct values) Roll a fair die $n$ times. Find the
expected number of faces that appear at least once. (Indicator for each face.)
\item (Cycle structure) In a uniformly random permutation of $n$ elements, show
that the expected number of cycles is $H_n$. (Hint: build the permutation by
the ``fundamental bijection'' or compute the expected number of cycles of
length $k$.)
\item Show that $\Var(X)\le\Exp[(X-c)^2]$ for every constant $c$, and use this
to show that if $a\le X\le b$ then $\Var(X)\le(b-a)^2/4$.
\item Let $X$ take values in $\{1,2,\dots\}$ with $\Prob(X=k)=\frac{6}{\pi^2k^2}$.
Show that $\Exp[X]=\infty$ but $\Prob(X<\infty)=1$.
\item (Challenge) Prove that the indicators $R_1,\dots,R_n$ of
Example~\ref{ex:records} are mutually independent, and conclude that the
variance of the number of records is $\sum_k\frac1k(1-\frac1k)$.
\end{enumerate}
